\documentclass[UTF8]{ctexart}
\usepackage{xeCJK}
\usepackage{geometry}
\geometry{left=2cm,right=2cm,top=2.5cm,bottom=2.5cm}
\setlength{\headheight}{12.7pt}

\usepackage{titlesec}
\titleformat{\section}
  {\normalfont\large\bfseries}
  {\thesection}
  {1em}
  {}
\titlespacing*{\section}
  {0pt}
  {3.5ex plus 1ex minus .2ex}
  {2.3ex plus .2ex}

\usepackage{amsmath}
\usepackage{amssymb}
\usepackage{amsthm}
\usepackage{enumitem}
\setlist[enumerate]{itemsep=0pt,topsep=0pt,parsep=0pt,leftmargin=0pt,itemindent=2.5em}
\setlist[itemize]{itemsep=0pt,topsep=0pt,parsep=0pt,leftmargin=0pt,itemindent=2.5em}
\usepackage{fancyhdr}
\pagestyle{fancy}
\fancyhf{}
\usepackage{graphicx}
\usepackage{hyperref}
\usepackage{indentfirst}
\usepackage{lmodern}

\begin{document}
\setlength{\abovedisplayskip}{1pt}
\setlength{\belowdisplayskip}{1pt}

\section{映射}

\hypertarget{sec:定义1.1}{}
\noindent\textbf{定义1.1 (映射)}

给定\href{02 二元关系#nameddest=sec:定义1.1}{二元关系} $f$, 如果
\[
    \forall x\,\forall y\,\forall y':\quad(x,y)\in f\land(x,y')\in f\to y=y',
\]
那么称$f$是一个\textbf{映射}.

此时对任意的$x\in\mathrm{Dom}\,f$, 存在唯一的$y$使得$(x,y)\in f$, 记作$f(x)$或$fx$.

\vspace{10pt}

给定映射$f$和集合$a,b$.
如果$\mathrm{Dom}\,f=a$且$\mathrm{Ran}\,f\subset b$, 就称$f$是从$a$到$b$的映射,
记作$f:a\to b$.

如果$\mathrm{Dom}\,f\subset a$且$\mathrm{Ran}\,f\subset b$,
就称$f$是从$a$到$b$的部分映射.

\vspace{10pt}

\hypertarget{sec:定义1.2}{}
\noindent\textbf{定义1.2 (映射集)}

给定集合$a,b$, 定义从$a$到$b$的\textbf{映射集}
\[
    ^ab\overset{\mathrm{def}}{=}\{f\subset a\times b\mid f:a\to b\},
\]
显然对任意的$f$, $f\in\,\!^ab$当且仅当$f$是从$a$到$b$的映射.

\vspace{10pt}

\hypertarget{sec:定理1.1}{}
\noindent\textbf{定理1.1 (映射的限制与复合)}

\begin{enumerate}
    \item 任给映射$f$和集合$a$, 限制$f\!\restriction_a$也是映射,
    并且对任意的$x\in\mathrm{Dom}\,f\!\restriction_a$, 有
    \[
        f\!\restriction_a(x)=f(x).
    \]
    \item 任给映射$f,g$, 复合$g\circ f$也是映射,
    并且对任意的$x\in\mathrm{Dom}(g\circ f)$, 有
    \[
        (g\circ f)(x)=g(f(x)).
    \]
\end{enumerate}

\noindent{\kaishu 证明.}

\begin{enumerate}
    \item 对任意的$x,y,y'$, 如果$(x,y),(x,y')$都属于$f\!\restriction_a$,
    那么它们也都属于$f$, 因此$y=y'=f(x)$.
    \item 对任意的$x,z,z'$, 如果$(x,z),(x,z')$都属于$g\circ f$, 则存在$y,y'$使得
    \[
        (x,y)\in f\land(y,z)\in g\quad\land\quad
        (x,y')\in f\land(y',z')\in g,
    \]
    于是$y=y'=f(x)$, 进一步$z=z'=g(f(x))$. \hfill $\qedsymbol$
\end{enumerate}

\vspace{10pt}

\hypertarget{sec:定理1.2}{}
\noindent\textbf{定理1.2}

给定映射$f$和集合$a,b$,
\begin{enumerate}
    \item 对任意的$x$, $x\in f^{-1}[a]$当且仅当$f(x)\in a$,
    \item $f^{-1}[a\cap b]=f^{-1}[a]\cap f^{-1}[b]$,
    \item $f^{-1}[a\setminus b]=f^{-1}[a]\setminus f^{-1}[b]$,
    \item $f[f^{-1}[a]]=a\cap\mathrm{Ran}\,f$.
\end{enumerate}





\newpage

\section{单射、满射与双射}

\hypertarget{sec:定义2.1}{}
\noindent\textbf{定义2.1 (单射、逆映射)}

给定映射$f$, 如果$f^{-1}$也是映射, 即
\[
    \forall x\,\forall x'\,\forall y:\quad(x,y)\in f\land(x',y)\in f\to x=x',
\]
那么称$f$是\textbf{单射}. 此时称$f^{-1}$是它的\textbf{逆映射}, 显然它也是单射.

\vspace{10pt}

\hypertarget{sec:定理2.1}{}
\noindent\textbf{定理2.1}

\begin{enumerate}
    \item 任给单射$f$和集合$a$, 限制$f\!\restriction_a$也是单射.
    \item 任给单射$f,g$, 复合$g\circ f$也是单射.
\end{enumerate}

\noindent{\kaishu 证明.}

\begin{enumerate}
    \item 对任意的$x,x',y$, 如果$(x,y),(x',y)$都属于$f\!\restriction_a$,
    那么它们也都属于$f$, 因此$x=x'$.
    \item $f^{-1},g^{-1}$都是映射, 所以$(g\circ f)^{-1}=f^{-1}\circ g^{-1}$
    也是映射. \hfill $\qedsymbol$
\end{enumerate}

\vspace{10pt}

\hypertarget{sec:定义2.2}{}
\noindent\textbf{定义2.2 (满射、双射)}

给定映射$f:a\to b$. 如果$\mathrm{Ran}\,f=b$, 就称$f:a\to b$是\textbf{满射}.

如果$f:a\to b$既是单射又是满射, 就称$f:a\to b$是\textbf{双射}.

\vspace{10pt}

\hypertarget{sec:定理2.2}{}
\noindent\textbf{定理2.2}

\begin{enumerate}
    \item 任给满射$f:a\to b$和$g:b\to c$, 复合$g\circ f:a\to c$也是满射.
    \item 任给双射$f:a\to b$和$g:b\to c$, 复合$g\circ f:a\to c$也是双射.
    \item 任给双射$f:a\to b$, 逆映射$f^{-1}:b\to a$也是双射.
\end{enumerate}

\noindent{\kaishu 证明.}

\begin{enumerate}
    \item $\mathrm{Ran}(g\circ f)=g[\,\mathrm{Ran}\,f\,]=g[b]=c$.
    \item 综上即得.
    \item 显然$f^{-1}$是单射, 并且$\mathrm{Dom}\,f^{-1}=\mathrm{Ran}\,f=b$,
    $\mathrm{Ran}\,f^{-1}=\mathrm{Dom}\,f=a$. \hfill $\qedsymbol$
\end{enumerate}

\vspace{10pt}

\hypertarget{sec:定义2.3}{}
\noindent\textbf{定义2.3 (恒等映射)}

对任意集合$a$, 记$\mathrm{id}_a\overset{\mathrm{def}}{=}\{(x,x)\mid x\in a\}$,
称作$a$上的\textbf{恒等映射}. 显然$\mathrm{id}_a$是$a$上的双射.

\vspace{10pt}

\hypertarget{sec:定理2.3}{}
\noindent\textbf{定理2.3}

给定集合$a,b$和映射$f:a\to b$, $g:b\to a$.
\begin{enumerate}
    \item 如果$g\circ f=\mathrm{id}_a$, 那么$f^{-1}\subset g$,
    并且$f$是单射, $g$是满射.
    \item 如果$f\circ g=\mathrm{id}_b$, 那么$f^{-1}\supset g$,
    并且$f$是满射, $g$是单射.
    \item 如果$g\circ f=\mathrm{id}_a$并且$f\circ g=\mathrm{id}_b$, 那么
    $f,g$都是双射且互为逆映射.
\end{enumerate}

\noindent{\kaishu 证明.}

\begin{enumerate}
    \item 对任意的$(y,x)\in f^{-1}$, 有$y=f(x)\implies x=g(f(x))=g(y)$,
    即$(y,x)\in g$. 所以$f^{-1}\subset g$.

    \hspace{2.17em}
    因为$g$是映射, 所以$f^{-1}\subset g$也是映射, 即$f$是单射.
    同时$\mathrm{Ran}\,g\supset\mathrm{Ran}\,f^{-1}=a$, 即$g$是满射.

    \item 由第一条即得.
    
    \item 由前两条即得. \hfill $\qedsymbol$
\end{enumerate}





\newpage

\section{选择公理}

\hypertarget{sec:定义3.1}{}
\noindent\textbf{定义3.1 (选择函数)}

任给集合$a$和映射$\displaystyle f:a\to\bigcup a$, 如果
\[
    \forall x\in a:\quad f(x)\in x,
\]
就称$f$是$a$上的一个\textbf{选择函数}.

\vspace{10pt}

\href{../01 ZFC 公理/01 ZFC 公理#nameddest=sec:选择公理}{选择公理}要求,
每个不含空集的集合都有选择函数. 据此有以下的定理.

\vspace{10pt}

\hypertarget{sec:定理3.1}{}
\noindent\textbf{定理3.1 (单值化定理)}

任给二元关系$r$, 存在映射$f$使得$f\subset r$并且$\mathrm{Dom}\,f=\mathrm{Dom}\,r$.

\noindent{\kaishu 证明.}

取$\mathcal{P}(\mathrm{Ran}\,r)\setminus\{\varnothing\}$上的选择函数$\varphi$.
对任意的$x\in\mathrm{Dom}\,r$, 定义
\[
    f(x)=\varphi(r[\{x\}]).
\]
此时对任意的$(x,y)\in f$, 有$y=\varphi(r[\{x\}])\in r[\{x\}]$, 即$(x,y)\in r$.
所以$f\subset r$. \hfill $\qedsymbol$

\vspace{10pt}

\hypertarget{sec:定理3.2}{}
\noindent\textbf{定理3.2 (单射定理)}

任给映射$f$, 存在单射$g$使得$g\subset f$并且$\mathrm{Ran}\,g=\mathrm{Ran}\,f$.

\noindent{\kaishu 证明.}

对$f^{-1}$应用单值化定理, 存在映射$h$使得$h\subset f^{-1}$并且
$\mathrm{Dom}\,h=\mathrm{Dom}\,f^{-1}=\mathrm{Ran}\,f$.

再令$g=h^{-1}$, 显然$g$是单射, $g\subset f$并且$\mathrm{Ran}\,g=\mathrm{Ran}\,f$.
\hfill $\qedsymbol$

\end{document}
